% ===
\section{Economic Commensalism}
% ===

* Definition. * Theoretical improvements. * Practical contributions. * Opportunities. * Carbon markets and public policies enhancements. * Economic and ecological consequences. * Economic Commensalism Index (ECI).

\subsection{Commensal Justification}

* Endogenization of the environment (\(E\)): Heuristic demonstration. * Pareto and Nash. * Cobb-Douglas, Arrow-Debreau and Fundamental Theorems of Welfare. * Isoquant and Isocost. * Tangency conditions (Lagrangean solution points). * Marginal effects on \(K\) and \(L\) (\(T\) and \(H\) also?). * Factor substitutability. * Effects on Solow Residual. * Extractivism, predation, mutualism, and commensalism at their optimum points. * Economic and ecological consequences.

\subsection{Heuristic Demonstration: A Microeconomic Reformulation}

* Modelagem heurística para mostrar a viabilidade do comensalismo. * Inserção de E na função de produção (ex.: Cobb-Douglas modificada). * Análise de tangência, substituibilidade, produtividade marginal. * Comparação com equilíbrios de Pareto e Nash. * Ilustração gráfica das curvas de produção sob diferentes regimes.

\subsection{The Economic Commensalism Index (ECI)}

* Proposta de um indicador para medir o grau de comensalismo em sistemas econômicos.

\subsection{Theoretical and Practical Implications}

* Como o comensalismo supera as limitações dos modelos atuais? * Implicações para políticas públicas, mercados de carbono, direito ambiental. * Transição para uma governança baseada em incentivos e não em coerção.

% fim %